\section{Algorithm}\label{sec:solution}

\subsection{Overview}
\label{subsec:algorithm-overview}

The algorithm follows the model introduced in Section~\ref{sec:model}.
It first privately selects a small set of useful predicate views, then constructs explanation candidates by pairing selected views with attributes, and finally privately selects and releases the top-$k$ histogram explanations.
This structure mirrors recent differentially private explanation frameworks, which separate the definition of explanation quality from the design of low-sensitivity private selection mechanisms \cite{TaoGMR22}.

Designing a practical algorithm for private query-difference explanations requires addressing two central challenges.
First, generating all possible views is infeasible because the space of predicate-defined subpopulations is combinatorial.
We therefore work with a candidate set $\mathcal{P}$ of predicate-based views, as is common in explanation and subgroup-discovery frameworks \cite{FEDEX,ICDE23}.
Second, the ideal utility function from Section~\ref{sec:explanation-quality} has high sensitivity and may be expensive to optimize.
We therefore introduce low-sensitivity proxy measures that capture the same conceptual goals: Histogram Distance for interestingness, Usefulness for representativeness and coverage, and DiversityScore for diversity.

Three refinements distinguish the algorithm presented here, each motivated by a specific failure of the direct formulation and each validated experimentally in Section~\ref{sec:experiments}:
\begin{enumerate}
  \item \textbf{Smoothed Histogram Distance} (Section~\ref{subsubsec:histogram-distance}). The unsmoothed score is a ratio whose global sensitivity is dictated by arbitrarily small histogram bins, even though almost all candidates rest on large ones. Smoothing the denominator by a public constant $\gamma$ yields a provably tighter bound and is the single change responsible for most of the empirical improvement.
  \item \textbf{Graded diversity} (Section~\ref{subsubsec:diversity-score}), which penalises the $j$-th reuse of an attribute by $\lambda^{j-1}$ rather than a flat $\lambda$, together with an order-independent set-level objective that makes Definition~\ref{def:setobj} well posed and monotone submodular.
  \item \textbf{Budget-derived $\tau$ and $\gamma$} (Section~\ref{subsubsec:adaptive-params}), following the annealing strategy of AIM~\cite{McKennaMSM22}: both are functions of the \emph{public} privacy accountant, so adapting them costs no privacy budget.
\end{enumerate}

\begin{figure}[ht]
\centering
\includegraphics[width=0.8\linewidth]{figures/placeholder.png}
\caption{An illustration of the private query-difference explanation framework.}
\label{fig:query-refinement}
\end{figure}
\alon{TODO: Replace the placeholder with the final framework figure.}

\subsection{Low-Sensitivity Proxy Measures}
\label{subsec:proxy-measures}

The proxy measures defined below are low-sensitivity counterparts of the ideal explanation-quality criteria introduced in Section~\ref{sec:explanation-quality}.
The View Usefulness score is used for private view selection; the Histogram Distance and DiversityScore are used for private explanation selection.

\subsubsection{View Usefulness}
\label{subsubsec:view-usefulness}

Coverage favors views that contain enough evidence from the compared query outputs:
\begin{equation}
Coverage(p,Q_1,Q_2,D)=\frac{|Q_1(p(D))|+|Q_2(p(D))|}{|D|}.
\label{eq:coverage-proxy}
\end{equation}
Representativeness favors views whose attribute distributions are not too far from the distribution of the full database:
\begin{equation}
Representativeness(p,D)=1-TVD(P_D,P_{p(D)}).
\label{eq:representativeness-proxy}
\end{equation}
\begin{equation}
Usefulness(p,Q_1,Q_2,D)=Representativeness(p,D)\cdot Coverage(p,Q_1,Q_2,D).
\label{eq:usefulness-proxy}
\end{equation}

Two conventions are fixed here that the earlier draft left open.
First, $Representativeness$ averages $1-TVD$ over a fixed set of attributes, not over the candidate attribute $A$: Usefulness is evaluated during \emph{view} selection, which precedes the pairing of views with attributes, so $A$ is not yet defined at that point.
Second, the average excludes the attribute(s) that the predicate $p$ itself constrains. Comparing $p(D)$ to $D$ on the very attribute that defines $p$ is close to tautological --- a view $\texttt{occupation}=c$ necessarily has an extreme occupation distribution --- and including it systematically suppresses exactly the single-attribute views that yield the clearest explanations.

\begin{proposition}[Sensitivity of View Usefulness]
\label{prop:usefulness-sensitivity}
Treating $|D|$ as public, $Usefulness$ has global sensitivity at most $2/|D|$.
\end{proposition}

\begin{proof}
$Representativeness\in[0,1]$ and $Coverage\in[0,2]$, and adding one tuple changes $|Q_1(p(D))|+|Q_2(p(D))|$ by at most $2$, hence changes $Coverage$ by at most $2/|D|$, while $Representativeness$ is a bounded multiplier.
The assumption that $|D|$ is public is necessary and is stated explicitly: $Coverage$ is a ratio whose denominator is itself data-dependent, so without it no constant bound independent of the database size exists.
\end{proof}

\subsubsection{Histogram Distance}
\label{subsubsec:histogram-distance}

Histogram Distance serves as a proxy for interestingness: an explanation is interesting when the distribution of an attribute changes substantially between $Q_1(p(D))$ and $Q_2(p(D))$.

\begin{definition}[Smoothed Histogram Distance]
\label{def:hd}
Given a dataset $D$, a view predicate $p$, an attribute $A$, two queries $Q_1,Q_2$, and a public smoothing constant $\gamma \ge 0$,
\begin{equation}
HD_\gamma(D;p,A,Q_1,Q_2)
=\!\!\sum_{v\in \operatorname{Dom}(D[A])}\!\!
\frac{\left|H_{Q_1(p(D)),A}(v)-H_{Q_2(p(D)),A}(v)\right|}{H_{p(D),A}(v)+\gamma}.
\label{eq:hd-definition}
\end{equation}
\end{definition}

Setting $\gamma=0$ recovers the unsmoothed score, in which case the sum must be restricted to bins with $H_{p(D),A}(v)>0$ to avoid division by zero.
For $\gamma>0$ no such restriction is needed: a bin absent from $p(D)$ is necessarily absent from $Q_1(p(D))$ and $Q_2(p(D))$ as well, since both are subsets of $p(D)$, so it contributes exactly $0/\gamma=0$.

The motivation for $\gamma$ is a mismatch between the \emph{worst case} the sensitivity bound must cover and the \emph{typical case} the mechanism actually encounters.
A bin containing two tuples can have its ratio swung by a full unit by a single individual, so the unsmoothed bound must accommodate that; yet the bound is then applied uniformly, including to candidates supported by thousands of tuples.
The exponential mechanism consequently cannot separate candidates whose scores differ by far less than the bound.
Smoothing raises every denominator by a fixed public amount, which barely perturbs well-supported bins while removing the pathological ones from consideration.

\begin{proposition}[Sensitivity of Smoothed Histogram Distance]
\label{prop:hd-sensitivity}
$GS(HD_\gamma)\le \dfrac{2}{1+\gamma}$.
\end{proposition}

\begin{proof}
Let $D'=D\cup\{t\}$.
Adding $t$ affects exactly one bin $v'$ of attribute $A$.
Write $f(a,b,m)=|a-b|/(m+\gamma)$ for the contribution of a bin with $Q_1$-count $a$, $Q_2$-count $b$ and view-count $m$.
For $v\neq v'$ all three quantities are unchanged, so those terms cancel.
For $v'$, each of $a,b,m$ changes by at most $1$. By the triangle inequality,
\[
|f(a,b,m)-f(a',b',m{+}1)| \le |f(a,b,m)-f(a,b,m{+}1)| + |f(a,b,m{+}1)-f(a',b',m{+}1)|.
\]
The first term is $|a-b|\left|\tfrac{1}{m+\gamma}-\tfrac{1}{m+1+\gamma}\right| \le \tfrac{1}{1+\gamma}$ since $|a-b|\le m$; the second is at most $\tfrac{|(a-b)-(a'-b')|}{m+1+\gamma}\le \tfrac{1}{1+\gamma}$.
Summing gives $GS(HD_\gamma)\le 2/(1+\gamma)$.
\end{proof}

The bound is strictly decreasing in $\gamma$, which is the entire point: $\gamma$ buys a tighter worst case at the cost of a mild, uniform compression of every candidate's score.
Section~\ref{subsec:gamma-tradeoff} reports this trade-off empirically, including the regime in which $\gamma$ becomes counter-productive.

\subsubsection{Diversity Score}
\label{subsubsec:diversity-score}

DiversityScore discourages repeatedly selecting explanations that describe the same aspect of the query difference.
Rather than a flat penalty, we grade the penalty by how many times an attribute has already been used:
\begin{equation}
DiversityScore(A,S)=\lambda^{\,\mathrm{count}(A,S)},
\qquad 0\le\lambda\le1,
\label{eq:diversity-score}
\end{equation}
where $\mathrm{count}(A,S)$ is the number of previously selected explanations already using attribute~$A$.
Setting $\lambda=0$ forbids repeats outright; $\lambda=1$ disables the diversity term; and the flat rule of the earlier draft is the special case in which every repeat beyond the first is charged the same $\lambda$.
The graded form matters only once an attribute can appear three or more times, hence only for larger $k$ --- with $k=3$ the two coincide almost surely, whereas at $k=10$ they differ in roughly half of all released sets.

Since $DiversityScore$ reads only previously selected \emph{outputs} and never the database, it adds no data sensitivity to the selection utility.

\subsubsection{Set-Level Objective}
\label{subsubsec:set-objective}

Equation~\eqref{eq:diversity-score} is defined relative to the explanations chosen in earlier iterations, so as written it scores an \emph{ordered sequence}, not a set --- and Definition~\ref{def:problem} asks for a set.
We therefore state the objective the algorithm approximates directly.

\begin{definition}[Set-level explanation utility]
\label{def:setobj}
For a set $E$ of candidates with scores $HD_\gamma(c)$ and attributes $A_c$,
\begin{equation}
F_\lambda(E)=\lambda\!\!\sum_{c\in E}\!HD_\gamma(c)
\;+\;(1-\lambda)\!\!\sum_{A}\max_{\substack{c\in E\\ A_c=A}}\!HD_\gamma(c).
\label{eq:set-objective}
\end{equation}
\end{definition}

Within each attribute group the strongest explanation counts in full and every further one counts at weight $\lambda$; more generally, assigning the $j$-th member of a group a weight $w_j$ with $w_1\ge w_2\ge\cdots\ge 0$ recovers Equation~\eqref{eq:diversity-score} for $w_j=\lambda^{j-1}$.

\begin{proposition}[Equivalence and structure]
\label{prop:setobj}
$F_\lambda(E)$ equals the total that the sequential objective~\eqref{eq:diversity-score} awards $E$ under its best ordering, and upper-bounds it under any ordering.
Moreover $F_\lambda$ is monotone submodular for any non-increasing weight profile.
\end{proposition}

\begin{proof}
For the first claim, the weights are non-increasing, so by the rearrangement inequality the ordering maximising $\sum_j w_j HD_{(j)}$ within a group is the one visiting that group in decreasing $HD_\gamma$, which is precisely what~\eqref{eq:set-objective} evaluates.
For the second, $\sum_{c\in E}HD_\gamma(c)$ is modular; each group term $\max_{c\in E,A_c=A}HD_\gamma(c)$ is monotone submodular because the marginal gain of adding $c$ is $\max(0,HD_\gamma(c)-\text{current group max})$, which cannot increase as $E$ grows; and a non-negative combination of monotone submodular functions is monotone submodular.
\end{proof}

Proposition~\ref{prop:setobj} sharpens Theorem~\ref{thm:non-private-hardness}.
Maximising a monotone submodular function subject to $|E|\le k$ is NP-hard, and the canonical reduction is from Maximum Coverage --- a sum of group maxima is exactly a weighted coverage function.
But the same structure guarantees that the greedy selection used by Algorithm~\ref{alg:topk} is a $(1-1/e)$-approximation, so the hardness result does not leave the algorithm without a guarantee.

\subsubsection{Budget-Derived $\tau$ and $\gamma$}
\label{subsubsec:adaptive-params}

Both $\tau$ and $\gamma$ admit natural data-dependent choices, but selecting either privately would consume budget that the selection and release stages need.
We instead follow the budget-annealing strategy of AIM~\cite{McKennaMSM22}, in which hyper-parameters are functions of the \emph{privacy accountant} rather than of the data:
\begin{equation}
\tau_t=\tau_{\min}+\tfrac{\epsilon_{\text{used}}}{\epsilon_S}\left(\tau_{\max}-\tau_{\min}\right),
\qquad
\gamma_t=\gamma_{\text{ref}}\cdot\frac{\epsilon_{\text{ref}}}{\epsilon_t}.
\label{eq:adaptive-params}
\end{equation}
Both depend only on public quantities, so by post-processing (Theorem~\ref{thm:post-processing}) they incur no privacy cost.
The form of $\gamma_t$ follows from Proposition~\ref{prop:hd-sensitivity}: the exponential mechanism's ability to separate candidates scales as $\epsilon_t/(2\,GS)=\epsilon_t(1+\gamma)/4$, so holding separation roughly constant as $\epsilon_t$ varies requires $\gamma_t\propto 1/\epsilon_t$ --- smooth harder when a round is noisy, less when it is sharp.

The reference value $\gamma_{\text{ref}}$ must still be fixed in advance from public information.
Section~\ref{subsec:gamma-tradeoff} shows that anchoring it to a fraction of $|D|$ is defensible but crude, because what $\gamma$ competes against is the size of individual histogram \emph{bins}, whose distribution relative to $|D|$ varies sharply across datasets.

\subsubsection{Combined Proxy Objective}
\label{subsubsec:combined-proxy-objective}

For private explanation selection over candidates $c=(p,A)$ we use
\begin{equation}
\widetilde{U}_i(D;c)=HD_{\gamma_t}(D;p,A,Q_1,Q_2)\cdot DiversityScore(A,S_{i-1}).
\label{eq:selection-utility}
\end{equation}

\begin{proposition}[Sensitivity of the selection utility]
\label{prop:selection-utility-sensitivity}
$GS(\widetilde{U}_i)\le 2/(1+\gamma_t)$.
\end{proposition}

\begin{proof}
Condition on the history $S_{i-1}$ of previously released outputs, which is legitimate under adaptive composition.
Given that history, $DiversityScore(A,S_{i-1})$ is a constant in $[0,1]$ determined entirely by released values, so it acts as a bounded multiplier on the data-dependent factor and cannot increase its sensitivity.
The claim follows from Proposition~\ref{prop:hd-sensitivity}.
\end{proof}

\subsection{Hardness of Non-Private Proxy Explanation Selection}
\label{subsec:non-private-hardness}

\begin{theorem}[Hardness of non-private proxy explanation selection]
\label{thm:non-private-hardness}
Selecting the optimal set of $k$ explanation candidates under the proxy objective is NP-hard, even when the proxy scores are known exactly and no privacy constraint is imposed.
\end{theorem}

\begin{proof}[Proof sketch]
By Proposition~\ref{prop:setobj} the objective is a weighted coverage function; Maximum Coverage reduces to it by giving each element unit score and each set an attribute group.
\alon{TODO: write the reduction out in full.}
\end{proof}

By Proposition~\ref{prop:setobj} the greedy procedure is nevertheless a $(1-1/e)$-approximation, and Section~\ref{subsec:quality} reports that in practice it matches the exhaustive optimum far more closely than that bound requires.

\subsection{Private Query-Difference Explanation Algorithm}
\label{subsec:private-algorithm}

\begin{algorithm}[t]
\caption{Private-Queries-Diff-Top-$k$-Explanation}
\label{alg:main}
\begin{algorithmic}[1]
\Require Private dataset $D$; queries $Q_1,Q_2$; views $\mathcal{P}$; attributes $Attrs$; view budget $\tau_{\max}$; explanations $k$; reference smoothing $\gamma_{\text{ref}}$; privacy budget $\epsilon$
\Ensure $k$ differentially private explanation histograms
\State $\epsilon_P,\epsilon_S,\epsilon_H \gets 0.25\epsilon,\ 0.45\epsilon,\ 0.30\epsilon$
\State \textbf{/* Private view selection */}
\State $u_{use}(p)\gets Usefulness(p,D,Q_1,Q_2)$
\State $\widehat{\mathcal{P}}\gets \textsc{OneShotTop-}\tau(\mathcal{P},u_{use}(\cdot),\epsilon_P,\Delta{=}2/|D|)$ \Comment{$\epsilon_P$-DP}
\Statex \textbf{(Post-processing)} Candidate construction incurs no privacy loss.
\State $\mathcal{E}\gets \textsc{Find-Top-$k$-Explanations}(D,Q_1,Q_2,\widehat{\mathcal{P}},Attrs,k,\gamma_{\text{ref}},\epsilon_S,\epsilon_H)$
\State \Return $\mathcal{E}$
\end{algorithmic}
\end{algorithm}

\begin{algorithm}[t]
\caption{Find-Top-$k$-Explanations}
\label{alg:topk}
\begin{algorithmic}[1]
\Require $D$; $Q_1,Q_2$; selected views $\widehat{\mathcal{P}}$; attributes $Attrs$; $k$; $\gamma_{\text{ref}}$; budgets $\epsilon_S,\epsilon_H$
\State $Counts\gets\emptyset$;\ \ $\mathcal{E}\gets\emptyset$;\ \ $\epsilon_t \gets \epsilon_S/k$
\For{$i\gets 1$ to $k$}
  \State $\tau_i \gets$ \Call{AdaptiveTau}{$\epsilon_{\text{used}},\epsilon_S,\tau_{\max}$} \Comment{public accountant}
  \State $\gamma_i \gets \gamma_{\text{ref}}\cdot \epsilon_{\text{ref}}/\epsilon_t$ \Comment{public; no privacy cost}
  \State $C_i \gets \{(p,A): p\in\widehat{\mathcal{P}}[1..\tau_i],\ A\in Attrs\}\setminus \mathcal{E}$
  \State $u_i(p,A)\gets HD_{\gamma_i}(D;p,A,Q_1,Q_2)\cdot \lambda^{Counts[A]}$
  \State $c_i=(p_i,A_i)\gets \textsc{ExponentialMechanism}\!\left(C_i,u_i,\epsilon_t,\Delta{=}\tfrac{2}{1+\gamma_i}\right)$
  \State $\Delta H_i(v)\gets H_{Q_1(p_i(D)),A_i}(v)-H_{Q_2(p_i(D)),A_i}(v)$
  \State $\widetilde{\Delta H_i}\gets \Delta H_i+\textsc{Laplace}(k/\epsilon_H)$
  \State $\mathcal{E}\gets\mathcal{E}\cup\{(p_i,A_i,\widetilde{\Delta H_i})\}$;\ \ $Counts[A_i]{+}{+}$
\EndFor
\State \Return $\mathcal{E}$
\end{algorithmic}
\end{algorithm}

The $k$ rounds of Algorithm~\ref{alg:topk} are exactly the one-shot top-$k$ mechanism, which is equivalent to $k$ invocations of the exponential mechanism~\cite{TaoGMR22}, applied here with the smoothed sensitivity of Proposition~\ref{prop:hd-sensitivity}.

\subsection{Differential Privacy Analysis}
\label{subsec:dp-proof}

\begin{proposition}[Sensitivity of noisy difference histogram release]
\label{prop:histogram-release-sensitivity}
For a fixed candidate $(p,A)$, $\Delta H_{p,A}$ has $L_1$-sensitivity at most $1$.
\end{proposition}

\begin{proof}
One tuple joins at most one bin of $A$ and belongs to at most one of $Q_1(p(D))$, $Q_2(p(D))$.
The affected coordinate changes by $d_1-d_2$ with $d_1,d_2\in\{0,1\}$, so by at most $1$; all other coordinates are unchanged, giving $\|\Delta H(D)-\Delta H(D')\|_1\le 1$.
\end{proof}

Note this yields $\textsc{Laplace}(k/\epsilon_H)$ in Algorithm~\ref{alg:topk}, rather than the $\textsc{Laplace}(2k/\epsilon_H)$ of the earlier draft; the factor of two was a conservative over-estimate of the same quantity.

\begin{proposition}[Differential Privacy of Algorithm~\ref{alg:topk}]
\label{thm:alg2-dp}
Algorithm~\ref{alg:topk} satisfies $(\epsilon_S+\epsilon_H)$-differential privacy.
\end{proposition}

\begin{proof}
Each round invokes the exponential mechanism with budget $\epsilon_S/k$ on a utility of sensitivity $2/(1+\gamma_i)$ (Proposition~\ref{prop:selection-utility-sensitivity}), which is $(\epsilon_S/k)$-DP, and the Laplace mechanism with budget $\epsilon_H/k$ on a statistic of $L_1$-sensitivity $1$ (Proposition~\ref{prop:histogram-release-sensitivity}), which is $(\epsilon_H/k)$-DP.
The quantities $\tau_i$ and $\gamma_i$ are functions of the public accountant only and are therefore not privacy-relevant.
The rounds are adaptive, since $u_i$ depends on previously selected attributes, but adaptive sequential composition applies \cite{DBLP:series/synthesis/2016Li}, giving $k(\epsilon_S/k+\epsilon_H/k)=\epsilon_S+\epsilon_H$.
\end{proof}

\begin{proposition}[Differential Privacy of Algorithm~\ref{alg:main}]
\label{thm:alg1-dp}
Algorithm~\ref{alg:main} satisfies $\epsilon$-differential privacy.
\end{proposition}

\begin{proof}
View selection is $\epsilon_P$-DP by Proposition~\ref{prop:usefulness-sensitivity} and the guarantee of the one-shot top-$\tau$ mechanism~\cite{TaoGMR22}.
Candidate construction is post-processing.
Algorithm~\ref{alg:topk} is $(\epsilon_S+\epsilon_H)$-DP by Proposition~\ref{thm:alg2-dp}.
Sequential composition gives $\epsilon_P+\epsilon_S+\epsilon_H=\epsilon$.
\end{proof}

\begin{corollary}[End-to-End Differential Privacy]
\label{cor:end-to-end-dp}
The overall explanation pipeline satisfies $\epsilon$-differential privacy.
\end{corollary}

\subsection{Complexity Analysis}
\label{subsec:complexity-analysis}

Let $|\mathcal{P}|$ be the number of candidate views, $m=|Attrs|$, $b$ the maximum number of bins per attribute, and $n=|D|$.
Computing $Usefulness$ for every view costs $O(|\mathcal{P}|\,n)$, and the one-shot top-$\tau$ mechanism adds $O(\tau|\mathcal{P}|)$.
Each of the $k$ rounds scores at most $\tau m$ candidates and releases one histogram, giving $O(k\,\tau\,m\,b)$ once the per-candidate bin vectors are materialised, plus $O(\tau m\, n)$ to materialise them.
The total is
\[
O\!\left(|\mathcal{P}|\,n \;+\; \tau m\, n \;+\; k\,\tau\,m\,b\right).
\]
The two $O(n)$ terms are a one-off preprocessing cost independent of $\epsilon$ and $k$; only the third is paid per round.
Because $HD_\gamma$ depends on $\gamma$ solely through the denominator, caching the per-bin numerator and denominator vectors makes the score computable for \emph{any} $\gamma$ without revisiting the data --- which is what keeps the adaptive $\gamma_t$ of Equation~\eqref{eq:adaptive-params} affordable, and is reflected in the measured separation of build and per-round cost in Section~\ref{subsec:scalability}.
